% 20_white_dwarf_predictions.tex
% Full derivation chain for AVE gravitational predictions on white dwarfs
\chapter{White Dwarf Gravitational Predictions}
\label{ch:white_dwarf_predictions}


\begin{objectivebox}
\begin{itemize}
    \item Derive white dwarf gravitational predictions (saturation-corrected redshift and standing shear-wave eigenfrequencies) from the saturated Diamond lattice impedance metric; relate them to the deflection and frame-dragging channels treated in Ch~\ref{ch:general_relativity}.
    \item Cross-validate the shear-eigenfrequency formula's shape against GR at a black-hole boundary (1.7\% on the quasi-normal-mode frequency), keeping the white-dwarf surface-mode prediction as a forward prediction.
\end{itemize}
\end{objectivebox}

This chapter applies the six-step AVE methodology (cf.\ Living Reference \S``How to Apply AVE'')
to the gravitational environment of white dwarf stars.  Two distinct predictions emerge:
\begin{enumerate}
    \item A saturation correction to the spectral gravitational redshift.
    \item A standing shear wave cavity with eigenfrequencies in the LIGO band.
\end{enumerate}
The derivation chain traces each claim from the four fundamental axioms,
through the universal operators, to a numerical prediction compared against observation.

%% ─────────────────────────────────────────────────────────────────
\section{Motivation: Why White Dwarfs?}

White dwarfs occupy a unique regime in the AVE landscape.
Their surface gravitational strain
\begin{equation}
    \varepsilon_{11}(R) = \frac{7GM}{c^2 R} \sim 10^{-3}
\end{equation}
is orders of magnitude larger than Earth's ($\sim 10^{-9}$),
placing their saturation correction well above the Regime I--II boundary noise floor
while remaining perturbatively small ($\varepsilon_{11} \ll \sqrt{2\alpha} \approx 0.12$).
Furthermore, high-precision spectroscopic redshift measurements exist
for several nearby white dwarfs, enabling direct comparison.

%% ─────────────────────────────────────────────────────────────────
\section{Step 1: LC Analogs}

The relevant constitutive parameters for the vacuum lattice
in the gravitational field of a white dwarf are:

\begin{center}
\begin{tabularx}{\textwidth}{@{}llX@{}}
\toprule
\textbf{Component} & \textbf{$\mu$-analog (inertia)} & \textbf{$\varepsilon$-analog (compliance)} \\
\midrule
Vacuum (exterior) & $\mu_0 \cdot n(r)$ & $\varepsilon_0 \cdot n(r)$ \\
WD interior (degenerate) & $\mu_{\text{plasma}}$ & $\varepsilon_0(1 - \omega_p^2 / \omega^2)$ \\
\bottomrule
\end{tabularx}
\end{center}

Under Symmetric Gravity (derived from Axioms~1 and~4), both $\mu$ and $\varepsilon$ scale by $n(r)$,
preserving the impedance: $Z(r) = \sqrt{\mu'/\varepsilon'} = Z_0$.

The WD interior is electron-degenerate matter with
$n_e \approx 7.5 \times 10^{35}\;\text{m}^{-3}$ (Sirius~B),
giving a plasma frequency $\omega_p \approx 5 \times 10^{20}\;\text{rad/s}$.
Since $\omega_p$ vastly exceeds any gravitational wave frequency,
the interior is \textit{evanescent} for shear perturbations:
the WD surface acts as a near-perfect reflector ($\Gamma \approx -1$).

%% ─────────────────────────────────────────────────────────────────
\section{Step 2: Strain and Regime}

For each white dwarf, we compute the principal radial strain and surface gravity:

\begin{center}
\begin{tabularx}{\textwidth}{@{}lcccX@{}}
\toprule
\textbf{White Dwarf} & $M/M_\odot$ & $R$ [km] & $\varepsilon_{11}$ & \textbf{Regime} \\
\midrule
Sirius B          & 1.018 & 5800 & $1.81 \times 10^{-3}$ & I (deep) \\
40 Eridani B      & 0.573 & 9000 & $6.58 \times 10^{-4}$ & I (deep) \\
Procyon B         & 0.602 & 8600 & $7.24 \times 10^{-4}$ & I (deep) \\
Stein 2051 B      & 0.675 & 8000 & $8.72 \times 10^{-4}$ & I (deep) \\
GD 358 (ZZ Ceti)  & 0.610 & 8800 & $7.17 \times 10^{-4}$ & I (deep) \\
\bottomrule
\end{tabularx}
\end{center}

All are deep Regime~I ($\varepsilon_{11} \ll r_{I \to II} = \sqrt{2\alpha} \approx 0.121$).
The saturation factor is perturbatively close to unity:
\begin{equation}
    S(\varepsilon_{11}) = \sqrt{1 - \varepsilon_{11}^2} \approx 1 - \frac{\varepsilon_{11}^2}{2}
\end{equation}

%% ─────────────────────────────────────────────────────────────────
\section{Step 3: Universal Operators}

\subsection{Prediction A: Saturation Correction to Gravitational Redshift}

The local clock rate in AVE is
\begin{equation}
    \frac{\omega_{\text{local}}}{\omega_\infty} = \frac{1}{n(R) \cdot S(\varepsilon_{11})}
\end{equation}
where $n(R) = 1 + 2GM/(c^2 R)$ is the gravitational refractive index (derived from Axioms~1 and~4)
and $S$ is the saturation factor (Axiom~4).

The gravitational redshift is therefore
\begin{equation}
    z_{\text{AVE}} = \frac{1}{\sqrt{1 - 2GM/(c^2 R)} \cdot S(\varepsilon_{11})} - 1
\end{equation}

For Standard General Relativity, $S = 1$ and
\begin{equation}
    z_{\text{GR}} = \frac{1}{\sqrt{1 - 2GM/(c^2 R)}} - 1
\end{equation}

The AVE correction over GR is therefore
\begin{equation}
    \delta z = z_{\text{AVE}} - z_{\text{GR}} = z_{\text{GR}} \cdot \left(\frac{1}{S} - 1\right) \approx z_{\text{GR}} \cdot \frac{\varepsilon_{11}^2}{2}
\end{equation}

Since $\varepsilon_{11} = 7\phi$ (where $\phi = GM/(c^2 R)$ is the Newtonian potential),
the correction scales as $49\phi^2/2$, which is \textbf{12.25 times larger} than the
standard PPN second-order correction $2\phi^2$.
This amplification arises from the Machian stress boundary $T_{\max} = c^4/(7G)$.

\begin{resultbox}{Sirius B Redshift Comparison}
\begin{center}
\begin{tabularx}{\textwidth}{@{}lrX@{}}
\toprule
\textbf{Quantity} & \textbf{Value} & \textbf{Source} \\
\midrule
$v_{\text{obs}}$  & $80.65 \pm 0.77$ km/s & Joyce et al.\ (2018) \\
$v_{\text{GR}}$   & 77.75 km/s           & Exact Schwarzschild \\
$v_{\text{AVE}}$  & 77.80 km/s           & $v_{\text{GR}} / S$ \\
Residual (obs$-$GR)  & $+2.90$ km/s    & \\
Residual (obs$-$AVE) & $+2.85$ km/s    & \\
AVE correction       & $\sim 0.05$ km/s & $z_{\text{GR}} \cdot \varepsilon_{11}^2/2$ \\
\bottomrule
\end{tabularx}
\end{center}
\end{resultbox}

\noindent
The AVE correction is in the correct direction (upward), but the 2.9~km/s residual
is dominated by the 3--5\% uncertainty in the mass-radius relation,
not by missing physics.
The prediction is therefore correct in direction but not currently discriminating:
at present mass-radius precision the $\sim$0.05~km/s saturation correction sits far
below the measurement residual and does not yet distinguish AVE from GR.

\subsection{Prediction B: Standing Shear Wave Eigenfrequencies}

The white dwarf surface acts as a boundary for shear perturbations
of the vacuum lattice.  Applying the 5-step regime-boundary eigenvalue method:

\begin{enumerate}
    \item \textbf{Boundary:} WD surface at $R$ (shear reflector, $\Gamma \approx -1$).
    \item \textbf{Effective cavity radius:} $r_{\text{eff}} = R / (1 + \nu_{\text{vac}}) = 7R/9$,
          where $\nu_{\text{vac}} = 2/7$ is the vacuum Poisson ratio.
    \item \textbf{Eigenfrequency:} $f_\ell = \ell \cdot c / (2\pi \, r_{\text{eff}})$.
    \item \textbf{Quality factor:} $Q = \ell$.
    \item \textbf{Decay time:} $\tau = Q / (\pi f)$.
\end{enumerate}

\begin{resultbox}{WD Shear Eigenfrequencies (l = 2)}
\begin{center}
\begin{tabularx}{\textwidth}{@{}lccX@{}}
\toprule
\textbf{White Dwarf} & $f_2$ [Hz] & \textbf{LIGO} & \textbf{Einstein Tel.} \\
\midrule
Sirius B         & 21.15 & \checkmark & \checkmark \\
40 Eridani B     & 13.63 & \checkmark & \checkmark \\
Procyon B        & 14.27 & \checkmark & \checkmark \\
Stein 2051 B     & 15.34 & \checkmark & \checkmark \\
GD 358 (ZZ Ceti) & 13.94 & \checkmark & \checkmark \\
\bottomrule
\end{tabularx}
\end{center}
\end{resultbox}

\noindent
\textbf{Cross-check (formula shape only, not WD-boundary corroboration):}
Applying the same formula to a Schwarzschild black hole---using the boundary at
$r_{\text{sat}} = 7GM/c^2$ instead of the WD surface $R$---reproduces the GR
quasi-normal mode frequency to 1.7\% accuracy
($\omega \cdot M_{\text{geom}} = 0.3673$ vs.\ GR exact $0.3737$ for $\ell = 2$).
This is a cross-regime consistency check on the formula's shape using a different
boundary radius; it does not corroborate the white-dwarf surface boundary choice.

%% ─────────────────────────────────────────────────────────────────
\section{Step 4: Symmetry Cancellations}

% eq_gravity_derived.tex — Gravity as Derived Consequence of Ax 1 + Ax 4 (Symmetric Scaling)
% Single source of truth for the gravitational refractive index, Symmetric Gravity
% derivation, and alpha-invariance under gravitational strain. \input this wherever
% the gravity content is formally restated.
%
% PROVENANCE: This file documents content that historically lived in eq_axiom_3.tex
% (Scheme B framing of "Axiom 3 = Gravity"). Per axiom homologation 2026-04-27 +
% Vol 1 Ch 1's canonical Scheme A: Axiom 3 is the Minimum Reflection Principle
% (eq_axiom_3.tex; renamed 2026-05-16 from "Effective Action Principle" — the new
% name follows the substrate-observable, the boundary reflection |Γ|², rather
% than the interior Lagrangian density). The gravitational refractive index n(r)
% and Symmetric Gravity scaling are derived consequences of Axiom 1 (LC network) +
% Axiom 4 (dielectric saturation applied symmetrically to mu and epsilon). Newton's
% constant G is a calibration constant emerging from Machian thermodynamic
% equilibrium of the LC lattice (per Vol 1 Ch 1:14-21 and Vol 1 Ch 1 Section 3.2).
%
\begin{resultbox}{Macroscopic Gravity (Derived from Ax 1 + Ax 4 Symmetric Scaling)}
Newton's constant sets the Machian boundary impedance:
\begin{equation}
    G = \frac{\hbar c}{7\xi \cdot m_e^2}
    \label{eq:gravity_newton}
\end{equation}
where $\xi = 4\pi(R_H / \ell_{node})\alpha^{-2} \approx 8.15 \times 10^{43}$ is the Machian hierarchy coupling.
The gravitational refractive index is
\begin{equation}
    n(r) = 1 + \frac{2GM}{rc^2}
    \label{eq:gravity_refraction}
\end{equation}
with Symmetric Gravity (Axiom 4 dielectric saturation applied to both inductive and capacitive constituents):
$\mu_{eff} = \mu_0 \cdot n(r)$,\;
$\varepsilon_{eff} = \varepsilon_0 \cdot n(r)$,\;
$Z = Z_0$ (invariant).
\end{resultbox}

\noindent\textbf{Derived Consequence 1: $\alpha$ Invariance.}
Under Symmetric Gravity, both constitutive parameters scale by the same factor $n \cdot S$
(combining gravitational refraction with Axiom 4 saturation). The fine-structure constant is therefore
\textit{exactly invariant} under gravitational strain:
\begin{equation}
    \alpha = \frac{e^2}{4\pi\,\varepsilon_{local}\,\hbar\,c_{local}}
           = \frac{e^2}{4\pi\,(\varepsilon_0 n S)\,\hbar\,(c_0/(nS))}
           = \frac{e^2}{4\pi\,\varepsilon_0\,\hbar\,c_0}
           = \alpha_0
    \label{eq:alpha_invariance}
\end{equation}
Multi-species clock comparisons at different gravitational potentials predict a
\textbf{null result} for $\Delta\alpha/\alpha$, consistent with all current experimental bounds.

\noindent\textbf{Derived Consequence 2: Temporal vs.\ Spatial Lattice Compression.}
The principal radial strain $\varepsilon_{11} = 7GM/(c^2 r)$ compresses the lattice asymmetrically.
The refractive index decomposes into:
\begin{align}
    n_{temporal} &= 1 + \tfrac{2}{7}\,\varepsilon_{11} && (\text{bulk/coordinate-time propagation index, slope 2}) \nonumber \\
    n_{spatial} &= 1 + \tfrac{9}{7}\,\varepsilon_{11} && (\text{controls matter-wave parallax, C11})
    \label{eq:lattice_decomposition}
\end{align}
% Notation cleanup 2026-05-17: both indices carry the "1 +" DC unit for parallelism.
% Both reduce to n = 1 in flat vacuum (eps_11 -> 0). The asymmetric prior form
% (n_spatial without "1 +") was shorthand for the strain-contribution term only,
% which produced confusion in driver-script implementations. The parity violation
% Delta_n = n_spatial - n_temporal = (9/7 - 2/7) * eps_11 = eps_11 is unchanged
% by the notation cleanup.
The temporal index $n_{temporal}$ (slope 2) is the \emph{bulk/coordinate-time}
propagation index --- what a signal traversing the gradient accumulates (Shapiro-class
integrated delay), $\approx 1/g_{00}$. The genuine \emph{local} clock rate / gravitational
redshift is a \emph{distinct} (slope 1) quantity, the proper tick of a clock sitting at $r$:
$\sqrt{g_{00}} = \sqrt{S} \approx 1 - GM/(c^2 r)$, giving redshift
$z \approx GM/(c^2 r)$. The two are bridged by $z = (n_{temporal} - 1)/2$: a propagating
signal picks up $2\times$ the local clock effect.
The spatial component governs the matter-wave (electron de~Broglie) parallax measured by
C11-MACH-ZEHNDER ($n_s - n_t = \varepsilon_{11}$) and frame dragging.
% W1 walk-back 2026-06-05 (Class-C internal-coherence; basis: research/2026-06-05_gravity-ppn-coherence-result.md):
% the (9/7) spatial index does NOT control light deflection. The canonical light-deflection
% derivation (Ch 2 \S double_deflection -> 4GM/bc^2) uses the (2/7) TRANSVERSE Cosserat-shear
% index n_perp = 1 + (2/7)chi_vol, not n_spatial. The (9/7) index IS load-bearing for the
% C11 electron matter-wave parallax (249.64 rad). KEEP-BOTH: (9/7) retained; only the
% "light deflection" attribution corrected.
% W2 walk-back 2026-06-05 (Class-C internal-coherence; Grant bulk-vs-local adjudication;
% basis: research/2026-06-05_gravity-sign-frequency-modulation-result.md):
% the "controls clock rate, redshift" annotation on n_temporal conflated TWO distinct
% temporal quantities. n_temporal (slope 2) = BULK/coordinate-time propagation index
% (integrated Shapiro delay, ~1/g00). The LOCAL clock rate / redshift is slope 1:
% sqrt(g00) = sqrt(S) ~ 1 - GM/rc^2, z = GM/rc^2 (the c_shear clock). Bridge:
% z = (n_temporal - 1)/2. RELABEL ONLY: n_temporal value (2/7 eps_11) and all
% derivations unchanged; only the "clock rate, redshift" label disambiguated from the
% slope-1 local clock. KEEP-BOTH.


\noindent
Two critical cancellations apply:

\begin{enumerate}
    \item \textbf{$\alpha$ is exactly invariant.}
    Under Symmetric Gravity, $\varepsilon_{\text{local}} = \varepsilon_0 \, n \, S$
    and $c_{\text{local}} = c_0 / (n \, S)$.
    In the definition of $\alpha$, the product $\varepsilon \cdot c$ is independent of $n \cdot S$:
    \begin{equation}
        \alpha = \frac{e^2}{4\pi \varepsilon_{\text{local}} \hbar \, c_{\text{local}}}
               = \frac{e^2}{4\pi \varepsilon_0 \hbar \, c_0} = \alpha_0
    \end{equation}
    Multi-species clock comparison experiments predict a null result for $\Delta\alpha / \alpha$.

    \item \textbf{Shear modes are spatial, not temporal.}
    Standing shear waves modify $g_{rr}$ (spatial metric),
    not $g_{00}$ (temporal metric), and therefore do not affect spectral redshift.
    Their observational signature is gravitational wave ringdown,
    not spectral line shift.
\end{enumerate}

%% ─────────────────────────────────────────────────────────────────
\section{Step 5: Numerical Engine Validation}

All predictions use engine constants from \texttt{ave.core.constants}:
\begin{verbatim}
from ave.core.constants import G, C_0, NU_VAC, M_SUN, ALPHA
from ave.gravity import (principal_radial_strain, saturation_radius,
                         shear_modulus_factor, refractive_index)
# Canonical regime path (identical physics):
# from ave.regime_3_saturated.orbital_impedance import (
#     calculate_refractive_strain)
\end{verbatim}

Scripts:
\begin{itemize}
    \item \kbleaf{src/scripts/vol\_4\_engineering/white\_dwarf\_saturation\_redshift.py}
    \item \kbleaf{src/scripts/vol\_4\_engineering/wd\_shear\_eigenfrequency.py}
\end{itemize}

%% ─────────────────────────────────────────────────────────────────
\section{Step 6: Testability}

\begin{center}
\begin{tabularx}{\textwidth}{@{}lllX@{}}
\toprule
\textbf{Observable} & \textbf{Frequency} & \textbf{Detector} & \textbf{Status} \\
\midrule
WD spectral redshift    & Optical        & HST/JWST   & Current (0.96\% precision) \\
WD $\ell=2$ shear mode  & 10--20 Hz      & LIGO/ET    & WD merger ringdown \\
$\Delta\alpha/\alpha$   & Multi-species  & NIST       & Null confirmed \\
BH QNM cross-check      & 100--200 Hz    & LIGO       & Validated to 1.7\% \\
\bottomrule
\end{tabularx}
\end{center}

%% ─────────────────────────────────────────────────────────────────
\section{Conclusions}

\begin{enumerate}
    \item The AVE framework predicts a unique nonlinearity in the gravitational
    redshift: the saturation factor $S(\varepsilon_{11})$ provides a correction
    that is 12.25 times larger than the standard PPN second-order term,
    due to the Machian stress amplification $\varepsilon_{11} = 7\phi$.

    \item For Sirius~B, this correction is $\sim$0.05~km/s---currently
    within the mass-radius uncertainty ($\sim$2.9~km/s residual),
    but potentially resolvable with improved M-R constraints from JWST.

    \item White dwarf surfaces act as vacuum shear reflectors,
    supporting standing wave eigenfrequencies at $\sim$10--20~Hz (LIGO band).
    These are a fundamentally different mode family from interior
    g-modes and p-modes observed in ZZ~Ceti asteroseismology.

    \item The fine-structure constant $\alpha$ is \textit{exactly invariant}
    under Symmetric Gravity, predicting a null result for all
    multi-species clock comparison experiments.

    \item The lattice density decomposes into temporal ($2/7 \cdot \varepsilon_{11}$)
    and spatial ($9/7 \cdot \varepsilon_{11}$) components.
    Only the temporal component contributes to gravitational redshift;
    the spatial component governs the matter-wave (electron de~Broglie) parallax
    (C11-MACH-ZEHNDER) and frame dragging. Light deflection couples instead to the
    transverse Cosserat-shear index $n_\perp = 1 + (2/7)\chi_{vol}$
    (Ch~\ref{ch:general_relativity}, \S\ref{sec:double_deflection}).
\end{enumerate}

\begin{summarybox}
\begin{itemize}
    \item White dwarf gravitational predictions (redshift, light deflection, frame dragging) are situated within the saturated Diamond lattice impedance metric---the same framework used for nuclear binding and water properties; this chapter derives the saturation-corrected redshift and the standing shear-wave eigenfrequencies, with light deflection and frame dragging treated in Ch~\ref{ch:general_relativity}.
    \item The temporal strain component governs gravitational redshift; the spatial component governs matter-wave parallax (C11) and frame dragging, while light deflection couples to the transverse index $n_\perp = 1 + (2/7)\chi_{vol}$ (Ch~\ref{ch:general_relativity} double-deflection $\to 4GM/bc^2$).
    \item The same eigenfrequency formula, applied to a Schwarzschild black hole at $r_{\text{sat}} = 7GM/c^2$, reproduces the GR quasi-normal-mode frequency to within $1.7\%$---a cross-regime shape validation of the impedance formula. The white-dwarf surface-mode prediction itself ($\sim$10--20\,Hz) is a forward prediction awaiting a WD-merger ringdown observation.
\end{itemize}
\end{summarybox}

\begin{exercisebox}
\begin{enumerate}
    \item Derive the gravitational redshift of Sirius~B from the AVE saturated lattice metric and compare with the GR prediction.
    \item Explain why the temporal and spatial components of strain separate into independent observables (redshift vs.~deflection), and identify an experiment that could distinguish AVE from GR at white dwarf surfaces.
\end{enumerate}
\end{exercisebox}

